% =============================================================================
\section{Transformação bilinear e pré-deformação}\label{sec:bilinear}
% =============================================================================

\subsection*{Teoria mínima}

A transformação bilinear mapeia o semiplano esquerdo $s$ no \emph{interior} do
círculo unitário:
\[
  \boxed{\ s=\frac{2}{T_s}\frac{1-z^{-1}}{1+z^{-1}}\ }
  \qquad\Longleftrightarrow\qquad
  z=\frac{1+\frac{T_s}{2}s}{1-\frac{T_s}{2}s}.
\]
Consequências:
\begin{itemize}
  \item \textbf{Não há aliasing}: todo o eixo $j\Omega$ é mapeado
        \emph{uma única vez} no círculo unitário.
  \item \textbf{Há distorção (warping)} no eixo de frequência:
        \[
          \Omega=\frac{2}{T_s}\tan\!\Bigl(\frac{\omega}{2}\Bigr)
          \ \Longleftrightarrow\
          \omega=2\arctan\!\Bigl(\frac{\Omega T_s}{2}\Bigr),
        \]
        que comprime frequências altas ($\Omega\to\infty\Rightarrow\omega\to\pi$).
  \item \textbf{Pré-deformação (pre-warping)} para manter as frequências
        críticas exatas:
        \[
          \boxed{\ \Omega_p=\frac{2}{T_s}\tan\!\Bigl(\frac{\omega_p}{2}\Bigr)
          =2F_s\tan\!\Bigl(\pi\frac{f_p}{F_s}\Bigr)\ }
        \]
\end{itemize}

\paragraph{Roteiro de projeto (passa-baixas/passa-faixa/altas/notch).}
\begin{enumerate}
  \item Converta as especificações digitais ($\omega_p,\omega_s$, dB) em
        frequências analógicas \emph{pré-deformadas} com a fórmula acima.
  \item Projete o protótipo analógico (Butterworth, Chebyshev, elíptico) com
        $\Omega_p,\Omega_s$. Para Butterworth de ordem $N$,
        \[
          N\ \ge\
          \frac{\log_{10}\!\left(\dfrac{\varepsilon_s}{\varepsilon_p}\right)}
               {\log_{10}\!\left(\dfrac{\Omega_s}{\Omega_p}\right)},
          \qquad
          \varepsilon=\sqrt{10^{0{,}1A}-1},\ A\ \text{em dB}.
        \]
  \item Aplique a bilinear ($s=\frac{2}{T_s}\frac{1-z^{-1}}{1+z^{-1}}$) e
        normalize $a_0=1$.
  \item Implemente em cascata de seções de 2ª ordem (Capítulo~\ref{sec:estr}).
\end{enumerate}

\paragraph{Passa-faixa.} O protótipo passa-baixas equivalente usa a
transformação $\Omega\to\frac{\Omega^2-\Omega_0^2}{B\Omega}$, com
$\Omega_0=\sqrt{\Omega_{p1}\Omega_{p2}}$ (centro geométrico) e
$B=\Omega_{p2}-\Omega_{p1}$ (largura). A ordem final do filtro passa-faixa é
$2N$ (cada polo do protótipo gera um par conjugado).

\paragraph{Comparação direta com a invariância ao impulso.}
\begin{center}
\begin{tabular}{p{4.1cm}p{4.6cm}p{4.6cm}}
\toprule
 & Invariância ao impulso & Bilinear\\
\midrule
Regra & $h[n]=T_sh_a(nT_s)$ & $s=\frac{2}{T_s}\frac{1-z^{-1}}{1+z^{-1}}$\\
Polos & $z=e^{pT_s}$ (exato) & $z=\frac{1+pT_s/2}{1-pT_s/2}$\\
Zeros & não mapeiam & mapeiam (mesma regra)\\
Aliasing & \textbf{sim} (réplicas) & \textbf{não}\\
Eixo de $\omega$ & linear & distorcido (exige pre-warping)\\
Passa-altas & \textbf{falha} & funciona\\
Nulos de $H$ & preservados em $\omega$ & preservados em $\Omega$\\
\bottomrule
\end{tabular}
\end{center}

\subsection*{Exercícios resolvidos}

% ---------------------------------------------------------------- 16.1
\exer{\textbf{(Desafio 09)} Com $F_s=1000$ Hz, calcule as frequências
analógicas pré-deformadas de $5$, $20$, $100$, $150$ e $300$ Hz e compare com
os valores sem pré-deformação.}

\textbf{Solução.} Aplicando $\Omega=2F_s\tan(\pi f/F_s)$:
\begin{center}
\begin{tabular}{cccc}
\toprule
$f$ [Hz] & $\omega=2\pi f/F_s$ & $\Omega$ com pre-warp [rad/s] & $2\pi f$ (sem warp)\\
\midrule
$5$   & $0{,}01\pi$ & $31{,}418511$  & $31{,}415927$\\
$20$  & $0{,}04\pi$ & $125{,}829335$ & $125{,}663706$\\
$100$ & $0{,}20\pi$ & $649{,}839392$ & $628{,}318531$\\
$150$ & $0{,}30\pi$ & $1019{,}050899$& $942{,}477796$\\
$300$ & $0{,}60\pi$ & $2752{,}763841$& $1884{,}955592$\\
\bottomrule
\end{tabular}
\end{center}
Observe como o erro relativo cresce com a frequência: $0{,}008\%$ em
$5$ Hz, $3{,}4\%$ em $100$ Hz e $46\%$ em $300$ Hz. Fazer pré-deformação só é
``opcional'' quando todas as frequências críticas são muito menores que
$F_s$.

\resp{$\Omega_{p1}=125{,}8293$, $\Omega_{p2}=1019{,}0509$,
$\Omega_{s1}=31{,}4185$, $\Omega_{s2}=2752{,}7638$ rad/s (os valores de
$5$ e $300$ Hz são os da banda de rejeição).}

% ---------------------------------------------------------------- 16.2
\exer{Projete por bilinear um passa-baixas de 1ª ordem com $f_c=100$ Hz e
$F_s=1000$ Hz. Verifique $|H|$ em DC, em $100$ Hz e em Nyquist.}

\textbf{Solução.} Pré-deformação: $\Omega_p=2\cdot1000\tan(0{,}1\pi)
=649{,}839392$ rad/s, e o protótipo é $H_a(s)=\frac{\Omega_p}{s+\Omega_p}$.
Substituindo $s=2000\frac{1-z^{-1}}{1+z^{-1}}$:
\[
  H(z)=\frac{\Omega_p(1+z^{-1})}{(2000+\Omega_p)+(\Omega_p-2000)z^{-1}}
      =\frac{0{,}2452372753\left(1+z^{-1}\right)}{1-0{,}5095254495\,z^{-1}} .
\]
Verificações:
\[
  |H(e^{j0})|=1 \quad(0\ \text{dB}),\qquad
  |H|_{\,f=100\,\text{Hz}}=\frac{1}{\sqrt2}=0{,}7071067812\
  (-3{,}0103\ \text{dB}),
\]
\[
  |H|_{\,f=F_s/2}=0 \quad(\text{zero em }z=-1),
\]
ou seja, a frequência de corte caiu \emph{exatamente} em $-3$ dB — o
pré-warping fez seu trabalho.

\resp{$b=\{0{,}2452372753;\ 0{,}2452372753\}$, $a=\{1;\ -0{,}5095254495\}$;
$-3{,}0103$ dB em $100$ Hz e zero em Nyquist.}

% ---------------------------------------------------------------- 16.3
\exer{Determine a ordem do protótipo Butterworth para o passa-faixa do
Desafio 09: banda passante de $20$ a $150$ Hz (ripple $1$ dB), rejeição em
$f\le5$ e $f\ge300$ Hz com $\ge30$ dB, $F_s=1000$ Hz.}

\textbf{Solução.} Usando os valores pré-deformados do Exercício 16.1:
\[
  B_p=\frac{\Omega_{p2}-\Omega_{p1}}{\Omega_{p1}\Omega_{p2}}
      =\frac{1019{,}050899-125{,}829335}{125{,}829335\cdot1019{,}050899}
      =0{,}0069659672,
\]
\[
  B_s=\frac{\Omega_{s2}-\Omega_{s1}}{\Omega_{s1}\Omega_{s2}}
      =\frac{2752{,}763841-31{,}418511}{31{,}418511\cdot2752{,}763841}
      =0{,}0314650993 .
\]
Fatores de ondulação:
\[
  \varepsilon_p=\sqrt{10^{0{,}1\cdot1}-1}=0{,}508847,\qquad
  \varepsilon_s=\sqrt{10^{0{,}1\cdot30}-1}=31{,}606961 .
\]
Ordem (protótipo passa-baixas equivalente):
\[
  N\ \ge\
  \frac{\log_{10}(\varepsilon_s/\varepsilon_p)}
       {\log_{10}(B_s/B_p)}
  =\frac{\log_{10}(62{,}11)}{\log_{10}(4{,}516975)}
  =\frac{1{,}7931}{0{,}6548}=2{,}738340
  \ \Longrightarrow\ \boxed{N=3}.
\]
O filtro digital final tem ordem $\mathbf{2N=6}$ (três seções de 2ª ordem),
com centro geométrico $\Omega_0=\sqrt{\Omega_{p1}\Omega_{p2}}
=358{,}087275$ rad/s ($56{,}9914$ Hz) e largura $B=\Omega_{p2}-\Omega_{p1}
=893{,}2216$ rad/s.

\resp{$N=3$ para o protótipo passa-baixas $\Rightarrow$ filtro passa-faixa de
ordem $6$ (3 seções SOS); $\varepsilon_p=0{,}5088$,
$\varepsilon_s=31{,}60696$, $B_s/B_p=4{,}5170$.}
